\chapter{How to Deal with Natural Disasters}

\section*{Definition and Overview}
Natural disasters, including bushfires, floods, cyclones, storms, and heatwaves, are recurring hazards in Australia. Dealing with them effectively means preparing before an event, acting safely during it, and supporting recovery afterwards. Preparation includes understanding local risks, having an emergency plan and kit, and knowing warning systems. During an event, the priorities are safety, evacuation or shelter as advised, and looking after physical and mental health. Recovery can take months to years, and seeking support for physical injuries, financial loss, and psychological distress is an important part of healing.

\section*{Epidemiology}
Natural disasters affect thousands of Australians every year. Bushfires, floods, and cyclones cause deaths, injuries, displacement, and enormous property and economic loss. Heatwaves cause more deaths than any other natural hazard. Mental health impacts are widespread: a large proportion of affected people experience anxiety, depression, and post-traumatic stress in the months after a disaster, with children, first responders, and people with existing vulnerabilities at higher risk. Prepared communities consistently experience lower death and injury rates.

\section*{Aetiology and Pathophysiology}
\begin{itemize}
    \item \textbf{Bushfires:} intense heat, flames, and smoke cause burns, smoke inhalation, and death
    \item \textbf{Floods:} fast-moving water causes drowning, injury, and water-borne infection
    \item \textbf{Cyclones and storms:} wind, debris, and lightning cause injury and structural damage
    \item \textbf{Heatwaves:} heat causes heat illness and worsens chronic disease
    \item \textbf{Displacement:} evacuation and loss of home disrupt medication, care, and support networks
    \item \textbf{Psychological stress:} threat, loss, and uncertainty trigger acute stress responses; prolonged exposure leads to anxiety, depression, and PTSD
    \item \textbf{Pathophysiology of injury:} burns, trauma, drowning, smoke inhalation, and cardiac events
\end{itemize}

\section*{Clinical Features}
\begin{itemize}
    \item Injuries: burns, fractures, lacerations, and crush injuries
    \item Smoke inhalation: cough, breathlessness, and carbon monoxide poisoning
    \item Heat illness and dehydration
    \item Infections after floods and in crowded evacuation centres
    \item Loss of medications and disruption of chronic disease management
    \item Acute stress: shock, anxiety, difficulty sleeping, and irritability
    \item Longer-term: depression, anxiety, PTSD, and grief
\end{itemize}

\section*{Differential Diagnosis}
\begin{itemize}
    \item Injuries must be assessed for severity and complications
    \item Breathlessness may be smoke injury, infection, or worsening heart or lung disease
    \item Psychological symptoms may be normal stress or develop into mental health conditions
    \item Delayed mental health effects should be considered in the months after a disaster
    \item Access to care may be limited, making triage and prioritisation important
\end{itemize}

\section*{When to Seek Medical Care}
Seek medical care for injuries, breathing problems, infections, and loss of medications. Psychological support should be sought early for persistent distress, and children and vulnerable people need particular attention.

\textbf{Call 000 (or local emergency services) for:}
\begin{itemize}
    \item Life-threatening injury, severe bleeding, or burns
    \item Difficulty breathing or smoke inhalation
    \item Drowning or near-drowning
    \item Chest pain or collapse
    \item Signs of heat stroke or severe dehydration
    \item Mental health crises with risk of self-harm
\end{itemize}

\section*{Diagnosis and Investigations}
\begin{itemize}
    \item \textbf{Risk assessment:} know your local hazards and warning systems
    \item \textbf{Emergency plans:} household plans cover evacuation, communication, and essential needs
    \item \textbf{Injury assessment:} examination and imaging for trauma and burns
    \item \textbf{Health needs:} medication, chronic condition, and special needs planning
    \item \textbf{Mental health screening:} for distress and trauma responses
    \item \textbf{Recovery assessment:} ongoing review of physical and emotional health
\end{itemize}

\section*{Management}
\begin{itemize}
    \item \textbf{Before:} prepare a plan, emergency kit, and insurance; know warnings
    \item \textbf{During:} follow official advice, evacuate early when advised, and prioritise life over property
    \item \textbf{First aid:} treat injuries, cool burns, and manage smoke exposure
    \item \textbf{Health continuity:} obtain replacement medications and maintain chronic disease care
    \item \textbf{Infection prevention:} safe water, hygiene, and wound care after floods
    \item \textbf{Psychological first aid:} support, routine, and connection after the event
    \item \textbf{Recovery support:} counselling, community recovery centres, and financial assistance
    \item \textbf{Children:} reassurance, routine, and age-appropriate information
\end{itemize}

\section*{Complications}
\begin{itemize}
    \item Burns, trauma, and their long-term consequences
    \item Smoke inhalation and lung injury
    \item Infections from contaminated water and crowded conditions
    \item Loss of medications and worsening of chronic conditions
    \item Post-traumatic stress, depression, and anxiety
    \item Financial loss, unemployment, and housing instability
    \item Disrupted education and child development
\end{itemize}

\section*{Prognosis and Outcomes}
Most people affected by natural disasters recover fully, particularly with good preparation and support. Physical injuries heal with prompt care, and most psychological distress resolves over months with family and community support. A minority develop longer-term mental health conditions that respond well to treatment. Communities that prepare, respond together, and sustain recovery support achieve better outcomes. Resilience is built through planning, connection, and accessing help early.

\section*{Prevention and Self-Care}
\begin{itemize}
    \item Prepare an emergency plan and kit before disaster season
    \item Stay informed through official warning systems and apps
    \item Act early: evacuate when advised rather than waiting
    \item Protect health with immunisation, chronic disease management, and medication supplies
    \item Look after mental health: sleep, routine, connection, and limiting distressing media
    \item Seek help early for persistent distress
    \item Support children with reassurance and information
    \item Engage with community recovery services after an event
\end{itemize}

\section*{Sources and Further Reading}
The following reputable sources provide further information for healthcare students and patients seeking reliable, current advice about dealing with natural disasters.

\begin{itemize}
    \item Australian Institute for Disaster Resilience. \textit{Emergency planning information.} https://knowledge.aidr.org.au/
    \item Emergency management agencies and state services. \textit{Local warnings and advice.} https://www.emergency.vic.gov.au/
    \item Healthdirect. \textit{Disaster and emergency health advice.} https://www.healthdirect.gov.au/
    \item Beyond Blue. \textit{Support after natural disasters.} https://www.beyondblue.org.au/
    \item Red Cross Australia. \textit{Preparedness and recovery resources.} https://www.redcross.org.au/
    \item World Health Organization. \textit{Health in emergencies.} https://www.who.int/
\end{itemize}
